% manuscript-sha256: 39b240b2592290928123eb7f457c8e0ceacfb513a59042eba77418841a6bdb37
% Options for packages loaded elsewhere
\PassOptionsToPackage{unicode}{hyperref}
\PassOptionsToPackage{hyphens}{url}
\documentclass[
]{article}
\usepackage{xcolor}
\usepackage[margin=1in]{geometry}
\usepackage{amsmath,amssymb}
\setcounter{secnumdepth}{-\maxdimen} % remove section numbering
\usepackage{iftex}
\ifPDFTeX
  \usepackage[T1]{fontenc}
  \usepackage[utf8]{inputenc}
  \usepackage{textcomp} % provide euro and other symbols
\else % if luatex or xetex
  \usepackage{unicode-math} % this also loads fontspec
  \defaultfontfeatures{Scale=MatchLowercase}
  \defaultfontfeatures[\rmfamily]{Ligatures=TeX,Scale=1}
\fi
\usepackage{lmodern}
\ifPDFTeX\else
  % xetex/luatex font selection
  \setmainfont[]{Times New Roman}
  \setmonofont[]{Consolas}
  \setmathfont[]{Latin Modern Math}
\fi
% Use upquote if available, for straight quotes in verbatim environments
\IfFileExists{upquote.sty}{\usepackage{upquote}}{}
\IfFileExists{microtype.sty}{% use microtype if available
  \usepackage[]{microtype}
  \UseMicrotypeSet[protrusion]{basicmath} % disable protrusion for tt fonts
}{}
\makeatletter
\@ifundefined{KOMAClassName}{% if non-KOMA class
  \IfFileExists{parskip.sty}{%
    \usepackage{parskip}
  }{% else
    \setlength{\parindent}{0pt}
    \setlength{\parskip}{6pt plus 2pt minus 1pt}}
}{% if KOMA class
  \KOMAoptions{parskip=half}}
\makeatother
\setlength{\emergencystretch}{3em} % prevent overfull lines
\providecommand{\tightlist}{%
  \setlength{\itemsep}{0pt}\setlength{\parskip}{0pt}}
\usepackage{bookmark}
\IfFileExists{xurl.sty}{\usepackage{xurl}}{} % add URL line breaks if available
\urlstyle{same}
\hypersetup{
  hidelinks,
  pdfcreator={LaTeX via pandoc}}

\author{}
\date{}

\allowdisplaybreaks
\begin{document}

\begin{center}
{\LARGE\bfseries Conditional density lifts and a primal–dual alignment statistic\par}
Jeremy H. Carroll\par
2026-10-23\par
Version 2.0.5\par
\end{center}

\subsection{Abstract}\label{abstract}

A nonzero real vector can be mapped to a unit vector by taking signed
square roots of its normalized absolute coordinates. The associated
rank-one projector retains relative magnitudes and signs while
forgetting total mass and common sign. We apply this construction to an
optimal residual and an optimal dual witness for a finite-dimensional
minimum-distance problem in the sum norm. Their trace distance defines
an alignment statistic. Complementary slackness gives a strict upper
bound below one for every optimal pair, and characterizes the
zero-distance case. Successive coordinate minimization selects a unique
pair once a coordinate order is fixed. The construction is undefined at
a zero residual; a separate convention assigns the statistic the value
zero there. We give proofs, an explicit example, and a selection policy
that uses the statistic only after primary residual comparisons. The
accompanying numerical implementation and discrete Lean results have
distinct, stated scopes.

\subsection{1. Problem and relation to existing
methods}\label{problem-and-relation-to-existing-methods}

Let \(B:\mathbb R^n\to\mathbb R^m\) be a fixed linear map, with
\(m\ge1\), and let \(\omega\in\mathbb R^m\). Define

\[
\Phi(\omega)=\min_{a\in\mathbb R^n}\|\omega-Ba\|_1.
\]

For a graph, \(B\) may be its vertex-to-edge incidence map. The minimum
is a norm on the quotient by \(\operatorname{im}B\) and a seminorm on
the original vector space. Linear programming supplies optimality and
duality statements {[}1{]}. The finite-dimensional projector distance
used below is standard {[}2{]}. The contribution here is the specified
combination of these constructions and its primal--dual bound, rather
than a new duality or trace-distance theorem.

\subsection{2. The signed square-root
map}\label{the-signed-square-root-map}

For \(v\ne0\), put \(M(v)=\sum_i|v_i|\) and

\[
\psi(v)_i=\operatorname{sgn}(v_i)\sqrt{|v_i|/M(v)},\qquad
\rho(v)=\psi(v)\psi(v)^T.
\]

The sign at zero is immaterial because that coordinate of \(\psi\) is
zero. We use \(\operatorname{sgn}(0)=1\) in the implementation.

\textbf{Proposition 2.1.} The matrix \(\rho(v)\) is positive
semidefinite, has trace one and rank one, and satisfies
\(\rho(v)^2=\rho(v)\). For \(c>0\), \(\rho(cv)=\rho(v)=\rho(-v)\).

\textbf{Proof.} The squared Euclidean norm of \(\psi(v)\) is
\(\sum_i|v_i|/M(v)=1\). Thus \(x^T\rho(v)x=(x^T\psi(v))^2\ge0\), the
trace is one, and multiplication of the two outer products gives
\(\rho(v)^2=\psi(v)(\psi(v)^T\psi(v))\psi(v)^T=\rho(v)\). Positive
scaling cancels in the normalized weights; negation changes the unit
vector's common sign and leaves its outer product unchanged.

For real unit vectors \(u,v\), the trace norm gives

\[
\frac12\|uu^T-vv^T\|_{\mathrm{tr}}
 =\sqrt{1-|u^Tv|^2}.
\]

Indeed, outside their span the difference is zero. On their span it has
trace zero and eigenvalues \(\pm\sqrt{1-|u^Tv|^2}\), with the
coincident-ray case obtained directly. This is a finite matrix identity,
with no physical interpretation assumed.

\subsection{3. Selecting an optimal
pair}\label{selecting-an-optimal-pair}

The residual and dual optimal sets are

\[
H^*(\omega)=\{h\in\omega-\operatorname{im}B:\|h\|_1=\Phi(\omega)\},
\]

\[
W^*(\omega)=\{\varphi:\|\varphi\|_\infty\le1,\ B^T\varphi=0,
\ \varphi^T\omega=\Phi(\omega)\}.
\]

Both sets are nonempty and compact. For the primal, the affine residual
space is closed and its intersection with any sufficiently large norm
ball is compact; a continuous norm attains its minimum there. The dual
feasible set is compact, and finite-dimensional norm duality gives the
stated equality of optimal values {[}1{]}. An optimal potential need not
be unique, but its residual belongs to \(H^*\).

Fix a coordinate order \(1,\ldots,m\). On either optimal set, minimize
the first coordinate, then the second on the remaining minimizers, and
continue. Each stage is nonempty and compact. At the last stage all
coordinates are fixed, so the remaining set contains exactly one vector.
Denote the resulting pair by \((h^*,\varphi^*)\). It is unique relative
to the declared coordinate order; it is not invariant under an arbitrary
change of that convention.

Support-first minimization is not a substitute for this construction. On
two disjoint directed triangles, a residual class can have optimal set

\[
\{h\ge0:h_1+h_2+h_3=1,\ h_4+h_5+h_6=1\}.
\]

Its lexicographically first possible sorted support is \((1,2,3,4)\).
Within that support, the first coordinate has infimum zero without
attaining it: \(h=(t,(1-t)/2,(1-t)/2,1,0,0)\) for \(0<t<1\). The stratum
is not closed.

\subsection{4. Alignment and its bounds}\label{alignment-and-its-bounds}

When \(\Phi(\omega)>0\), define

\[
\Xi(\omega)=\frac12\|\rho(h^*)-\rho(\varphi^*)\|_{\mathrm{tr}}.
\]

When \(\Phi(\omega)=0\), define \(\Xi(\omega)=0\) separately. This value
is a convention: the density lift of the zero residual does not exist.

\textbf{Theorem 4.1.} For any optimal primal--dual pair \((h,\varphi)\)
with \(\Phi>0\) and \(N=\|\varphi\|_1\),

\[
1\le N\le m,\qquad
0\le\frac12\|\rho(h)-\rho(\varphi)\|_{\mathrm{tr}}
\le\sqrt{1-1/N}\le\sqrt{1-1/m}<1.
\]

\textbf{Proof.} Dual feasibility and optimality imply
\(\varphi^Th=\varphi^T\omega=\|h\|_1\). Each summand
\(|h_i|-\varphi_i h_i\) is nonnegative, so it is zero. Thus
\(\varphi_i=\operatorname{sgn}(h_i)\) whenever \(h_i\ne0\). In
particular \(\|\varphi\|_\infty=1\), whence \(1\le N\le m\).
Consequently

\[
\langle\psi(h),\psi(\varphi)\rangle
=\frac{\sum_{h_i\ne0}\sqrt{|h_i|}}{\sqrt{\Phi N}}
\ge\frac{\sqrt{\sum_i|h_i|}}{\sqrt{\Phi N}}
=\frac1{\sqrt N}.
\]

The trace-distance identity gives the result.

\textbf{Corollary 4.2.} The distance is zero exactly when the two
supports coincide and \(|h_i|=\Phi/N\) on that support. Their signs then
agree automatically.

\textbf{Proof.} Zero distance means equal rank-one projectors, hence
equal unit vectors up to a common sign. Their overlap is positive by
Theorem 4.1, so the common sign is positive. Comparing their squared
coordinates gives \(|h_i|/\Phi=|\varphi_i|/N\). On the primal support,
\(|\varphi_i|=1\); outside it equality requires \(\varphi_i=0\). The
converse follows by substitution.

\subsection{5. Worked example and a selection
policy}\label{worked-example-and-a-selection-policy}

On a directed triangle take \(h=(2,1,1)\) and \(\varphi=(1,1,1)\). Their
pairing and primal mass are both four, and the circulation constraint
\(B^T\varphi=0\) holds. These are an optimal pair for the class of
\(h\). Here \(N=3\) and

\[
\langle\psi(h),\psi(\varphi)\rangle=\frac{1+\sqrt2}{\sqrt6},\qquad
\Xi=\sqrt{\frac{3-2\sqrt2}{6}}>0.
\]

This pair illustrates why sign and support agreement alone do not imply
zero distance. It need not be the pair selected by coordinate
minimization.

One possible decision policy compares tuples \((\Phi,\Gamma,R,\Xi)\)
with a resolution \(\epsilon>0\). It accepts a candidate if one of these
conditions holds:

\begin{enumerate}
\def\labelenumi{\arabic{enumi}.}
\tightlist
\item
  \(\Phi'<\Phi-\epsilon\).
\item
  \(\Phi-\epsilon\le\Phi'\le\Phi\) and \(\Gamma'<\Gamma-\epsilon\).
\item
  \(\Phi-\epsilon\le\Phi'\le\Phi\),
  \(\Gamma-\epsilon\le\Gamma'\le\Gamma\), \(R'\le R_{\max}\) and
  \(\Xi'<\Xi-\epsilon\).
\end{enumerate}

This policy never increases the primary coordinate. The ratio condition
is applied only in the third branch; primary or secondary improvement
can be accepted without it. This is the specified behavior, rather than
a general feasibility or convergence result. All tuple entries must be
computed or validated for their stated mathematical roles before
interpreting a decision.

\subsection{6. Implementation and formal
scope}\label{implementation-and-formal-scope}

The Python module \nolinkurl{code/mod_conditional_density_lift_core.py}
implements the numerical square-root lift and trace distance. Its
numerical function accepts supplied vectors; it does not certify their
optimality or perform the successive coordinate selector. The
linear-program examples use floating-point optimization and cannot serve
as exact certificates merely because a solver reports success. The
explicit unverified-override option is outside the mathematical
definition.

An optional zero-obstruction certificate is bound to the supplied
cochain and checks its coboundary equality over the exact rational
values represented by the finite binary inputs. A numerical
near-equality does not certify an exact zero class. The numerical lift
refuses inputs whose total mass cannot be stored as a finite
floating-point number, rather than returning a clipped mass.

The module also contains exact integer/rational checks for normalized
absolute weights, the Cauchy inequality, complementary slackness and
rational decision comparisons. \nolinkurl{lean4/Xi.lean} formalizes
corresponding discrete statements. The real rank-one matrix
construction, compactness argument and square-root bound above are prose
proofs, not results of the discrete Lean development. The supplied
census lists the declarations audited by the pinned compiler.

\subsection{7. Limitations}\label{limitations}

The statistic depends on a selected optimal pair and a coordinate
convention. It discards total residual mass and cannot replace the
quotient residual as a measure of that mass. Its definition at zero is
separate from the density formula. A small statistic does not imply
feasibility, convergence, or preservation of constraints. The finite
matrix encoding does not derive quantum mechanics or any law of physical
measurement.

\subsection{Publication information}\label{publication-information}

Version DOI:
\href{https://doi.org/10.5281/zenodo.23604860}{10.5281/zenodo.23604860}.
Concept DOI:
\href{https://doi.org/10.5281/zenodo.20192296}{10.5281/zenodo.20192296}.
The public record identifies the October 23, 2026 edition as version
2.0.5.

\subsection{References}\label{references}

\begin{enumerate}
\def\labelenumi{\arabic{enumi}.}
\tightlist
\item
  S. Boyd and L. Vandenberghe, \emph{Convex Optimization}, Cambridge
  University Press, 2004, Chapters 4--5.
  \href{https://web.stanford.edu/~boyd/cvxbook/}{Author-hosted book}.
\item
  J. Watrous, \emph{The Theory of Quantum Information}, Cambridge
  University Press, 2018, Chapter 3.
  \href{https://cs.uwaterloo.ca/~watrous/TQI/}{Author-hosted text and
  chapters}.
\end{enumerate}

\end{document}
